
\chapter*{General Introduction  
%\markboth {Introduction }{}}
\markboth {General Introduction}{}}
\addcontentsline{toc}{chapter}{General Introduction}



% \lettrine[lines=2]{T}\\he automobile business is continually developing, fueled by technology advancements that are redefining cars. From engine control and advanced driver assistance systems, to infotainment systems, software is crucial in regulating the different components of a vehicle's operation. Nowadays automotive software engineering has gained at least as much attention and importance as any other engineering aspect of the automotive industry process.
% \newline
% \newline
% \indent
% With the modern era of interconnected systems and high-speed communication, the automotive industry had also merged the lines of this technological advancement. The vehicle of the future is getting shipped with advanced driver assistance systems (ADAS), connected vehicles, over-the-air (OTA) update, autonomous driving and recently with the shift towards electric vehicles (EVs) that further amplifies the importance of software, as it orchestrates battery management, energy efficiency, and charging infrastructure integration. With this increase of complexity and the need for precision as well as real time operations, these systems rely on accurate timing information for various functioning. Thus the need for precise time synchronization has become paramount and the demand for such a module has grown significantly.



\newacronym{adas}{ADAS}{Advanced Driver Assistance Systems}
\newacronym{evs}{EVs}{Electric Vehicles}
\newacronym{ota}{OTA}{Over-the-air}
\newacronym{v2x}{V2X}{Vehicle to Everything}
\newacronym{ecus}{ECUs}{Electronic Control Units}
\newacronym{ptp}{PTP}{Precision Time Protocol}
\newacronym{ieee}{IEEE}{Institute of Electrical and Electronics Engineers}


\lettrine[lines=2]{I}\\n the contemporary industrial landscape, characterized by fast-paced technological advancements and heightened market demands, automation and robotics have emerged as fundamental drivers of innovation and productivity. These technologies revolutionize manufacturing processes by enhancing precision, reliability, and speed while minimizing human error and operational costs. The integration of automated systems not only streamlines production but also empowers industries to adapt swiftly to changing requirements, thus maintaining a competitive edge in the global marketplace.


Against this backdrop, the present final year project was conducted within the premises of Xpert-Meca, a dynamic company renowned for its expertise in the design, development, and implementation of tailor-made automated and robotic solutions. Xpert-Meca’s commitment to advancing industrial automation makes it an ideal environment for undertaking projects that combine theoretical knowledge with practical application.


The project’s core objective was to conceive, develop, and simulate an automated cube-sorting cell. This system was specifically designed to demonstrate the potential and versatility of modern automation technologies in an industrial context, particularly as an exhibit for trade fairs and industrial exhibitions. The sorting cell not only illustrates the synergy between robotics and automation control but also serves as a practical showcase of how complex processes can be managed efficiently through integrated systems.

The work was systematically organized into four comprehensive chapters, each addressing distinct but interconnected aspects of the project:


\begin{itemize}
\item \textbf{Chapter 1} introduces the host company, Xpert-Meca, providing insight into its organizational structure, areas of expertise, and external collaborations. This chapter also outlines the complete project framework, including the problem statement, objectives, the chosen solution, and the high-level operating principle of the automated sorting cell;
\item \textbf{Chapter 2} details the system's design and architecture, covering the hardware and software components. It provides a technical overview of the operative parts (KUKA robot, conveyors, pneumatic systems, and sensors), the control parts (PLC, HMI, and robot controller), and the supporting electrical and pneumatic cabinets. This chapter also outlines the software tools used for control, programming, design, and visualization;
\item \textbf{Chapter 3} delves into the control system and automation of the robotic sorting cell, detailing its automation logic and programming. It explains how a Programmable Logic Controller (PLC) serves as the system's central brain, prioritizing safety and sequential operation. The chapter also covers the use of the GRAFCET graphical language for designing the control logic, and the TIA Portal software for a unified engineering framework that integrates hardware configuration, PLC programming, and HMI design. Finally, it details the PLC programming implementation, including hardware configuration, I/O and tag management, and the use of various programming languages like Ladder Diagram (LAD) and fundamental programming blocks;
\item \textbf{Chapter 4}  delves into the system integration and verification of the automated sorting cell. It details the PROFINET communication framework used to connect the PLC, robot, and other components via a central router. The chapter also explains the crucial role of virtual commissioning by demonstrating how software like TIA Portal and Blender were used to create a "digital twin." This allowed for comprehensive simulation and verification of the control logic, HMI functionality, and mechanical movements, ensuring a robust and reliable design before physical implementation.

\end{itemize}

Throughout the project, I tackled significant technical challenges, from initial design constraints and component compatibility to programming and system integration hurdles. The solutions I developed combine analytical thinking with practical engineering skills and adherence to best industry practices. Ultimately, this endeavor provided a valuable learning experience, bridging academic knowledge with industry requirements and offering a solid foundation for future work in industrial automation and robotics.